\chapter{Prise en main}

\section{Objet du manuel}

Ce manuel décrit l'utilisation quotidienne de ChargeTrackr pour les trois profils
qui travaillent directement avec le réseau de recharge :

\begin{longtable}{|L{3.2cm}|L{5.1cm}|L{7.1cm}|}
  \hline
  \rowcolor{ctmint}\textbf{Profil} & \textbf{Mission principale} & \textbf{Périmètre dans ChargeTrackr} \\ \hline
  \endfirsthead
  \hline
  \rowcolor{ctmint}\textbf{Profil} & \textbf{Mission principale} & \textbf{Périmètre dans ChargeTrackr} \\ \hline
  \endhead
  Client & Recharger un véhicule et suivre ses dépenses &
  Rechercher une station, choisir un connecteur, démarrer ou arrêter une session,
  consulter ses paiements, reçus et abonnements. \\ \hline
  Opérateur & Superviser l'activité du réseau &
  Contrôler les stations et connecteurs de son organisation, traiter les alertes,
  coordonner la maintenance, suivre les sessions et produire des rapports. \\ \hline
  Technicien & Exécuter les travaux de terrain &
  Consulter les stations, traiter les alertes et interventions assignées, joindre
  des preuves et remettre un rapport de maintenance. \\ \hline
\end{longtable}

Les écrans et données visibles dépendent du rôle et de l'organisation de
l'utilisateur. Un technicien ne dispose donc pas des mêmes actions qu'un
opérateur, même lorsqu'ils consultent la même station.

\section{Prérequis}

La version en ligne est accessible à l'adresse
\url{https://chargetrackr.me}. Pour utiliser l'application, il faut :

\begin{itemize}[leftmargin=1.6em]
  \item un navigateur récent avec JavaScript et les cookies activés ;
  \item une connexion au serveur ChargeTrackr ;
  \item un compte actif et une adresse électronique vérifiée ;
  \item l'autorisation de géolocalisation uniquement pour la recherche
  \emph{Near me} ;
  \item un accès à l'adresse électronique associée au compte pour l'activation
  ou la réinitialisation du mot de passe.
\end{itemize}

\attention{La position du client n'est demandée qu'au moment d'une recherche de
proximité. Le rayon utilisé est réglable dans la page \textbf{Settings}.}

\section{Créer un compte client}

Le client peut créer lui-même un compte depuis le lien
\textbf{Create a client account} de la page de connexion. Cette inscription
publique ne crée ni organisation, ni compte opérateur ou technicien.

\begin{procedurebox}{Inscription avec une adresse électronique}
  \begin{enumerate}[leftmargin=1.6em]
    \item ouvrir \textbf{Create a client account} ;
    \item saisir les informations demandées et choisir un mot de passe robuste ;
    \item valider le formulaire ;
    \item ouvrir le courriel de vérification envoyé par ChargeTrackr ;
    \item suivre le lien, puis se connecter.
  \end{enumerate}
\end{procedurebox}

Le client peut aussi utiliser Google. ChargeTrackr associe alors le compte à
l'identité Google sans enregistrer le mot de passe Google. Un compte créé
uniquement avec Google ne propose pas de changement de mot de passe local.

\section{Se connecter}

\begin{procedurebox}{Connexion avec une adresse électronique}
  \begin{enumerate}[leftmargin=1.6em]
    \item ouvrir la page \textbf{Sign in} ;
    \item saisir l'adresse électronique du compte ;
    \item saisir le mot de passe ;
    \item sélectionner \textbf{Sign in to dashboard} ;
    \item vérifier que l'espace correspondant au rôle est affiché.
  \end{enumerate}
\end{procedurebox}

\manualscreenshot{login.png}{Page de connexion à ChargeTrackr}{fig:user-login}

Si le mot de passe local est oublié, l'action \textbf{I forgot my password}
envoie un lien temporaire à l'adresse vérifiée. Le nouveau mot de passe n'est
enregistré qu'après validation de ce lien.

\expected{L'utilisateur arrive sur son tableau de bord. Le libellé situé près de
son avatar confirme le rôle actif.}

\section{Activer un compte invité}

Les comptes opérateur et technicien sont créés sous forme d'invitation par
l'administrateur de leur organisation. L'administrateur ne choisit jamais le mot
de passe à la place de l'employé.

\begin{procedurebox}{Activation depuis le courriel d'invitation}
  \begin{enumerate}[leftmargin=1.6em]
    \item ouvrir le courriel d'invitation ChargeTrackr ;
    \item suivre le lien d'activation avant son expiration ;
    \item vérifier l'identité, le rôle et l'organisation proposés ;
    \item définir puis confirmer un mot de passe robuste ;
    \item compléter les informations facultatives proposées ;
    \item valider et se connecter.
  \end{enumerate}
\end{procedurebox}

\securitynote{Ne jamais transférer le lien d'activation. S'il est expiré ou
révoqué, demander un nouvel envoi à l'administrateur de l'organisation.}

\section{Première connexion}

La page d'accueil personnalisée présente un parcours court orienté vers un
premier résultat utile. Elle ne remplace pas l'application : elle prépare le
contexte du compte, puis renvoie vers le véritable espace de travail.

\begin{itemize}[leftmargin=1.6em]
  \item le client découvre comment trouver une station et démarrer une recharge ;
  \item l'opérateur identifie les vues de supervision et les alertes ;
  \item le technicien repère ses tâches assignées et la remise de rapport.
\end{itemize}

Le parcours peut être quitté et repris ultérieurement depuis le menu du compte.

\expected{À la fin de l'accueil personnalisé, l'utilisateur entre dans son
espace réel. Les conseils contextuels peuvent ensuite apparaître au moment où
une action est exécutée, sans bloquer la navigation.}

\section{Se déconnecter}

Ouvrir le menu de l'avatar, puis choisir \textbf{Logout}. Fermer une fenêtre sans
se déconnecter n'est pas recommandé sur un poste partagé.
